%---------- 第三章 方法与实验 ----------
\section{改进方法与实验分析}

\subsection{改进网络结构设计}

\quad 针对复杂场景下目标尺度变化剧烈、背景干扰严重导致的识别精度下降问题，本文在
基线网络的基础上引入通道—空间双重注意力模块：通道注意力用于筛选显著性特征通道，
空间注意力用于定位目标区域。改进后的网络在骨干网络的第 3、4 个残差阶段之后分别
插入注意力模块，并以跳跃连接保留原始特征信息，避免注意力加权造成的信息损失。

\subsection{实验设置}

\quad 实验环境为 Ubuntu 22.04，GPU 为 NVIDIA RTX 4090，深度学习框架为 PyTorch 2.2。
训练采用 SGD 优化器，初始学习率为 0.1，动量为 0.9，权重衰减为 $5 \times 10^{-4}$，
批大小为 128，共训练 200 个轮次，并在第 60、120、160 轮将学习率衰减为原来的
$1/10$。

\subsection{实验结果与分析}

\quad 本文方法与基线方法在 CIFAR-10 测试集上的对比结果如表~\ref{tab-results} 所示。
可以看出，引入注意力机制后，识别准确率提升了 1.6 个百分点，且参数量仅增加 3.2\%，
表明改进方法以极小的额外开销获得了明显的性能增益。

\begin{table}[htbp]
  \centering
  \caption{不同方法在 CIFAR-10 测试集上的对比结果}
  \label{tab-results}
  \begin{tabular}{lccc}
    \toprule
    方法 & 准确率（\%） & 参数量（M） & 推理时延（ms） \\
    \midrule
    ResNet-18（基线）        & 94.8 & 11.2 & 6.4 \\
    SE-ResNet-18             & 95.5 & 11.8 & 7.0 \\
    本文方法                 & 96.4 & 11.6 & 6.9 \\
    \bottomrule
  \end{tabular}
\end{table}

\quad 此外，消融实验表明：单独引入通道注意力可将准确率提升 0.9 个百分点，单独引入
空间注意力可提升 1.1 个百分点，二者结合可获得最优效果，说明通道与空间注意力在
特征增强中具有互补作用。
